\chapter{State Objects}
\label{ch:state-objects}

CNA's four state-object classes all inherit \cnaclass{GraphicsResource}. Their static,
ready-made presets exist before a \cnaclass{GraphicsDevice}; that is why the base accepts a null
device (Chapter~\ref{ch:graphicsdevice}). An architectural choice worth stating up front is that
CNA stores the three device-wide state objects \emph{by value}. A device assignment copies the
supplied state into its own CPU field, then calls the corresponding backend application method.
Thus mutating the original local value later cannot change the stored device value. CNA also
declares the static presets \texttt{const}, whereas
FNA's \texttt{static readonly} fields protect only the C\# reference --- their public state
properties remain mutable. The presets themselves are unbound, harmless \cnaclass{GraphicsResource}
values in both implementations, not GPU allocations that must be disposed.

There is a second, more important rule than that value/reference difference: reading CNA's
mutable device-state getter gives access only to the CPU copy. Changing that reference does
\emph{not} call \texttt{ApplyBlendState}, \texttt{ApplyDepthStencilState}, or
\texttt{ApplyRasterizerState}; those calls occur only in their three device setters. Reassign a
changed copy to submit it. A former CNA test comment claimed that FNA's reference alias would make
a later \cnaclass{BlendState} mutation take effect on the next draw. Current FNA source does not
support that conclusion: its blend/depth flush resubmits only when the selected state \emph{object
identity} changes, so a mutation of the already-current object is not dirty-tracked. FNA does
reapply its rasterizer state every draw, but that implementation detail is not a portable update
route. The safe rule in either API is to configure a state before assignment and assign it again
after a change.

\section{A worked example: the by-value semantics, made concrete}

The value copy and submission boundary are easy to state abstractly and easy to get wrong in
practice. \texttt{GraphicsDevice::setBlendStateProperty(const BlendState\&)} takes its argument
by \texttt{const} reference, stores an internal copy, and submits the supported fields. The
mutable overload of \texttt{getBlendStateProperty()} returns that CPU copy; it is not a proxy that
submits every field mutation:

\begin{lstlisting}
BlendState custom;
custom.setColorSourceBlendProperty(Blend::SourceAlpha);
custom.setColorDestinationBlendProperty(Blend::One);   // additive-style blend
device.setBlendStateProperty(custom);

custom.setColorSourceBlendProperty(Blend::Zero);        // mutate the ORIGINAL afterward
// device's stored copy and submitted blend factors are still SourceAlpha here.

// Do not use this as a GPU update: it changes only the device's stored CPU copy.
device.getBlendStateProperty().setColorSourceBlendProperty(Blend::One);

BlendState submitted = device.getBlendStateProperty();
submitted.setColorSourceBlendProperty(Blend::One);
device.setBlendStateProperty(submitted); // the required submission step
\end{lstlisting}

The practical rule is therefore more conservative than the old ``mutate through the getter''
advice: construct or copy, change, and assign a fresh value every time the effective blend,
depth-stencil, or rasterizer state changes. A C++ getter mutation can make later introspection
misleading by showing a value the backend was never asked to use.

\section{The value surface is broader than the binding bridge}
\label{sec:state-object-binding-boundary}

Every public property round-trips in the small C\texttt{++} value objects, but the common backend
interface is narrower. This is a structural limitation above any individual backend mapping:
the three whole-state setters make one immediate call with the fields present in that interface;
there is no later generic reconciliation of omitted fields. The current routes are:

\begin{itemize}
  \item \cnaclass{BlendState} forwards its six source/destination/function values. Its
    \texttt{BlendFactor} is then sent through the separate device property. Each color-write
    mask and \texttt{MultiSampleMask} is absent from the backend call. The device's own
    \texttt{MultiSampleMask} is only a C\texttt{++} integer shadow; unlike FNA's native setter,
    it makes no backend call.
  \item \cnaclass{DepthStencilState} forwards all of its depth/stencil fields on whole-state
    assignment. Its \texttt{ReferenceStencil} is also sent through the distinct device setter,
    whose backend default is a no-op; Chapter~\ref{ch:state-objects}'s EasyGL caveat therefore
    still applies to a standalone reference change.
  \item \cnaclass{RasterizerState} forwards cull mode, fill mode, scissor enable, and both bias
    values, but not \texttt{MultiSampleAntiAlias}. The property is a retained CPU value rather
    than a generic rasterization command.
  \item On every CNA draw, the device visits all sixteen \emph{pixel} sampler slots and forwards
    only filter, AddressU, AddressV, and maximum anisotropy. AddressW, MaxMipLevel, and
    MipMapLevelOfDetailBias are never read by that bridge. The public vertex texture and vertex
    sampler collections are retained (and cleared when a texture is disposed), but no draw path
    forwards either collection to the backend at all.
\end{itemize}

The sampler loop is intentionally different from the three whole-state setters: changing a
pixel sampler collection entry is seen at the next draw because every slot is reapplied then.
That does not rescue its omitted fields or the vertex collections. Current CPU tests exercise
defaults, names, round trips, and bounds, while backend examples exercise selected supported
pixel-sampler and state fields. No located test asks whether a getter mutation reaches a backend,
or distinguishes ColorWriteChannels, either multisample property, AddressW, MaxMipLevel, LOD
bias, or a vertex sampler/texture from its default behavior.

\section{BlendState}

Four presets --- \texttt{Additive}, \texttt{AlphaBlend}, \texttt{NonPremultiplied},
\texttt{Opaque} --- and a property surface covering the color and alpha blend function/source/%
destination triples independently, four independent \texttt{ColorWriteChannels} masks (one
per potential MRT slot), a \texttt{BlendFactor} color, and a \texttt{MultiSampleMask}.
\cnaclass{Blend} has thirteen values (\texttt{One}, \texttt{Zero}, the four
source/destination color and alpha pairs, \texttt{BlendFactor} / \texttt{InverseBlendFactor},
and \texttt{SourceAlphaSaturation}); \cnaclass{BlendFunction} has five (\texttt{Add},
\texttt{Subtract}, \texttt{ReverseSubtract}, \texttt{Max}, \texttt{Min}).
\texttt{ColorWriteChannels} is a flags enum with the same hand-implemented bitwise operator
set already seen on \cnaclass{DisplayOrientation} (Chapter~\ref{ch:game-loop}) and
\cnaclass{ClearOptions} (Chapter~\ref{ch:graphicsdevice}).

For the six source/destination/function values that the common bridge actually carries, the
current cross-backend status in \texttt{docs/graphics-backend-feature-matrix.md} is fully
correct on EasyGL, BGFX, SDL\_Renderer, and all three Direct3D backends. Vulkan is the one
backend with a real, historically serious bug in that forwarded subset --- for a long time its
blend implementation was, in the project's own words, ``almost entirely fake'': one hardcoded
blend equation applied regardless of what was actually requested, confirmed failing across five
separate pixel tests before the fix replaced it with a genuine per-\cnaclass{Blend}/%
\cnaclass{BlendFunction} mapping across all nine of the backend's pipeline-creation sites. That
evidence does not make the unforwarded color-write masks or multisample mask work.

\section{DepthStencilState}

Three presets (\texttt{Default} --- depth test and write on, \texttt{LessEqual} ---,
\texttt{DepthRead} --- test on, write off ---, and \texttt{None} --- both off) sit atop a
sixteen-property surface: \texttt{DepthBufferEnable} / \texttt{DepthBufferWriteEnable}/%
\texttt{DepthBufferFunction}, and a full stencil-test configuration
(\texttt{StencilEnable}, \texttt{StencilFunction}, \texttt{StencilMask}/%
\texttt{StencilWriteMask}, \texttt{ReferenceStencil}, the three stencil-operation slots for
pass/fail/depth-fail, \texttt{TwoSidedStencilMode}, and four separate counter-clockwise-face
variants of the same operations). \cnaclass{CompareFunction} has eight values;
\cnaclass{StencilOperation} has eight (\texttt{Keep}, \texttt{Zero}, \texttt{Replace},
\texttt{Increment} / \texttt{Decrement} and their saturating variants, \texttt{Invert}).

Depth testing itself is solid: \texttt{DepthBufferWriteEnable} and the full five-value
\texttt{DepthBufferFunction} sweep both verify correctly on EasyGL and (as of a fix) Vulkan.
Stencil testing is where this state object's real history lives. At one point, Vulkan's
apply-state entry point accepted fifteen parameters but genuinely stored only two --- every
stencil-related parameter passed through was simply discarded, and \texttt{stencilTestEnable}
was never set on any pipeline at all. Five separate, independent checks (enable flag, masks,
front-face operations, two-sided mode, reference-stencil propagation) each separately
confirmed failing on Vulkan while passing on EasyGL, before the fix added real per-pipeline
compare operations, full front/back \texttt{VkStencilOpState} configuration, and dynamic-state
reference/mask values. A second Vulkan-specific bug was found in the same pass: depth-format
selection checked a stencil-\emph{less} format before it checked any stencil-capable one.
EasyGL, meanwhile, had a subtler bug of its own: no window was ever created requesting an
actual stencil buffer (\texttt{SDL\_GL\_STENCIL\_SIZE}) at all --- meaning every stencil test
that appeared to pass was passing against a buffer with zero stencil bits actually allocated,
a ``correct code, no resource to act on'' bug rather than a logic error. One gap named in an
earlier pass of this chapter as open on both EasyGL and BGFX is now only open on one of them:
\texttt{IGraphicsBackend::SetReferenceStencil} does exist in the shared backend interface, with
a no-op default any backend may override --- reading each backend's own header directly finds
BGFX now has a real override (Task~764, mirroring Vulkan's own existing mechanism, Task~870/%
319: every other stencil parameter is cached so a later standalone reference-value change can
rebuild the backend's own front/back stencil state without a full state-object
re-application), while EasyGL still has none, falling through to the interface's own no-op ---
independently overriding \texttt{ReferenceStencil} without reassigning the whole
\cnaclass{DepthStencilState} is confirmed broken on EasyGL specifically, not BGFX anymore. A
second gap named in the same earlier pass of this chapter as still open on all three hardware
backends is also closed now: a stencil-only \texttt{Clear} call once never even checked its
own flag on any backend (Task~871), but a real fix now dispatches it
correctly on EasyGL, Vulkan, and BGFX alike, surfacing and closing a genuine, independent
Vulkan-specific bug along the way (every render-pass-creation site had its own stencil
load-operation hardcoded to discard, regardless of what \texttt{Clear()} itself requested) ---
see Chapter~\ref{ch:graphicsdevice} for the fix in full.

\subsection{A worked example: masking a reflection with the stencil buffer}

The classic real use of the stencil-test properties above, taken together, is a planar mirror:
render the mirror surface into the stencil buffer only, then render the reflected scene a
second time, but only where that stencil mark was left behind --- so the reflection never
spills past the mirror's own on-screen footprint even though the reflected geometry itself may
extend well beyond it. This needs two distinct \cnaclass{DepthStencilState} configurations, not
one:

\begin{lstlisting}
DepthStencilState writeMirrorMask;
writeMirrorMask.setDepthBufferEnableProperty(true);
writeMirrorMask.setStencilEnableProperty(true);
writeMirrorMask.setStencilFunctionProperty(CompareFunction::Always);
writeMirrorMask.setStencilPassProperty(StencilOperation::Replace);
writeMirrorMask.setReferenceStencilProperty(1);

DepthStencilState maskedByMirror;
maskedByMirror.setDepthBufferEnableProperty(true);
maskedByMirror.setStencilEnableProperty(true);
maskedByMirror.setStencilFunctionProperty(CompareFunction::Equal);
maskedByMirror.setReferenceStencilProperty(1); // only draw where the mirror already wrote 1

// Pass 1: draw the mirror quad itself, writing 1 into the stencil buffer everywhere it covers.
device.setDepthStencilStateProperty(writeMirrorMask);
mirrorQuad.Draw();

// Pass 2: draw the reflected scene, masked to the mirror's own footprint by the stencil test.
device.setDepthStencilStateProperty(maskedByMirror);
reflectedScene.Draw(); // a full scene draw, using a reflected view matrix
\end{lstlisting}

One of this chapter's own already-stated facts directly constrains whether this exact
technique is safe to ship across all three hardware backends. \texttt{ReferenceStencil} above
is set once, as part of constructing each \cnaclass{DepthStencilState} --- this worked example
does not rely on changing it independently afterward, which sidesteps the still-open gap
flagged above (no real \texttt{SetReferenceStencil} override exists on EasyGL specifically,
BGFX's own fix now notwithstanding; a technique needing a per-draw-call reference value without
rebuilding the whole state object would hit that remaining gap directly on EasyGL). Task~871
made the requested stencil \emph{value} real on all three hardware backends, but did not make a
stencil-only clear selective everywhere. EasyGL issues only the stencil bit; Vulkan/BGFX clear
all pass/view attachments at pass start and can therefore disturb colour/depth too
(\S\ref{sec:backend-clear-matrix}). Portable code should keep relying on what this technique's
own first pass already does --- its unconditional
\texttt{CompareFunction::Always} / \texttt{StencilOperation::Replace} configuration overwrites
the stencil value outright every frame, never depending on it having been zeroed first, which
is exactly why this specific technique remains correct without assuming selective clear
semantics.

\section{RasterizerState}

Three presets (\texttt{CullClockwise}, \texttt{CullCounterClockwise} --- the XNA default ---
\texttt{CullNone}) and six properties: \texttt{CullMode}, \texttt{DepthBias},
\texttt{FillMode}, \texttt{MultiSampleAntiAlias}, \texttt{ScissorTestEnable},
\texttt{SlopeScaleDepthBias}. \cnaclass{CullMode} has three values
(\texttt{None}, \texttt{CullClockwiseFace}, \texttt{CullCounterClockwiseFace});
\cnaclass{FillMode} has two (\texttt{Solid}, \texttt{WireFrame}).

\texttt{ScissorTestEnable} is deliberately only an enable switch. Its rectangle lives
separately in \texttt{GraphicsDevice.ScissorRectangle}; changing either property immediately
routes only that half of the state to the backend. The distinction is observable: EasyGL,
Vulkan, BGFX, WebGPU, SDL GPU, D3D9, and D3D11 retain an independent Boolean and rectangle.
The SDL renderer and ASCII instead apply the rectangle as an always-active clip and ignore the
Boolean. D3D12, Software, Canvas, and DX3 never turn the pair into a custom raster clip at all;
Headless only validates and traces it. The full scope and evidence matrix is
\S\ref{sec:backend-viewport-scissor-matrix}. Also note the shared
\cnaclass{SpriteBatch} routing gap from Chapter~\ref{ch:spritebatch}: passing this state to
\texttt{SpriteBatch::Begin()} currently discards it before any backend sees it.

All three cull modes verify correctly on both EasyGL and Vulkan via a genuine
opposite-winding-order pixel test, and a small architectural detail worth noting: Vulkan
flips clip-space Y and \emph{compensates} by setting \texttt{VK\_FRONT\_FACE\_CLOCKWISE}
instead of the API's default counter-clockwise convention --- confirmed, by the same test, to
cull the same input winding identically to EasyGL despite the two backends taking different
routes to get there. \texttt{FillMode::WireFrame} is verified on both EasyGL (via a
\texttt{GL\_LINES} re-expansion, since OpenGL ES has no \texttt{glPolygonMode} to fall back
on directly) and Vulkan (\texttt{VK\_POLYGON\_MODE\_LINE}). \texttt{DepthBias} and
\texttt{SlopeScaleDepthBias} have one known, pre-existing failure at an extreme magnitude
(\texttt{DepthBias = -1e6}) on Vulkan specifically --- the other three sub-cases pass, and
this was not re-investigated as part of this state object's own audit. BGFX's
\texttt{CullMode} / \texttt{FillMode} / \texttt{DepthBias} / \texttt{ScissorTestEnable} are, by
design, not independently pixel-verified at all --- BGFX has no GPU-readback API in this
project, so its state coverage is smoke-test/no-regression only, a scope difference worth
remembering whenever this book or the project's own docs describe a BGFX row as
``correct.''

\subsection{A worked example: shadow-map depth bias, and a wireframe debug toggle}

\texttt{DepthBias} and \texttt{SlopeScaleDepthBias} exist for exactly one recurring real
problem: shadow acne, the self-shadowing Moir\'e artifact a shadow-mapped surface shows when its
own depth-buffer sample is compared against itself with no tolerance at all. A small, fixed bias
plus a slope-scaled term (larger at grazing angles, where acne is worst) is the standard fix:

\begin{lstlisting}
RasterizerState shadowMapBias;
shadowMapBias.setCullModeProperty(CullMode::CullClockwiseFace);
shadowMapBias.setDepthBiasProperty(0.0015f);
shadowMapBias.setSlopeScaleDepthBiasProperty(1.5f);

device.setRasterizerStateProperty(shadowMapBias);
// ... render the scene from the light's own view/projection into the shadow map ...
\end{lstlisting}

This chapter's own audit notes one specific, narrow failure at an extreme magnitude
(\texttt{DepthBias = -1e6}) on Vulkan --- ordinary shadow-mapping bias values like \texttt{0.0015}
above are nowhere near that magnitude and are unaffected; the caveat only matters if a game
computes its bias value dynamically and that computation could, under some real input, blow up
toward an extreme.

A second, unrelated everyday use for the same class --- toggling a wireframe debug overlay ---
is a one-property change with no bias involved at all:

\begin{lstlisting}
RasterizerState wireframe;
wireframe.setFillModeProperty(FillMode::WireFrame);
wireframe.setCullModeProperty(CullMode::None); // see both front and back edges while debugging

if (showWireframeOverlay)
    device.setRasterizerStateProperty(wireframe);
else
    device.setRasterizerStateProperty(RasterizerState::CullCounterClockwise);
\end{lstlisting}

Remember this chapter's own scope caveat before trusting a wireframe-toggle test on BGFX
specifically: \texttt{FillMode} there has no GPU-readback pixel verification at all (BGFX has no
readback API in this project), so \texttt{FillMode::WireFrame} on that one backend is
smoke-tested, not pixel-proven, unlike the same toggle on EasyGL or Vulkan.

\section{SamplerState}

Six presets --- \texttt{AnisotropicClamp} / \texttt{Wrap}, \texttt{LinearClamp} / \texttt{Wrap},
\texttt{PointClamp} / \texttt{Wrap} --- cover exactly the combinations real FNA ships, no more.
Seven properties: \texttt{AddressU} / \texttt{V} / \texttt{W} (\cnaclass{TextureAddressMode}:
\texttt{Wrap}, \texttt{Clamp}, \texttt{Mirror}), \texttt{Filter} (\cnaclass{TextureFilter},
nine values spanning point/linear/anisotropic and their per-mip-stage combinations),
\texttt{MaxAnisotropy}, \texttt{MaxMipLevel}, \texttt{MipMapLevelOfDetailBias}.

The central, now-fixed finding here was not about the address-mode or filter mapping tables
themselves --- both were correct in isolation for as long as they were tested --- but about
whether the supported \emph{pixel} \texttt{GraphicsDevice.SamplerStates} route was actually
applied before a 3D draw. It was not, for three separate, backend-specific reasons. EasyGL's
eighteen user-primitive overloads omitted the internal sampler-application step; only the
ordinary buffer-bound \texttt{DrawPrimitives} and \texttt{DrawIndexedPrimitives} calls made it.
On Vulkan, the dual-texture descriptor-set path hardcoded a default sampler into its slots
regardless of what sampler state was actually
current; on BGFX, a dual-texture draw bound its second texture slot using the \emph{first}
slot's sampler flags. All three were fixed, and --- once actually applied --- the underlying
address-mode and filter mappings themselves were confirmed correct on EasyGL and Vulkan for a
real 3D stock-effect draw. That conclusion never extended to \texttt{VertexSamplerStates}:
the public collection is still retained only, as the binding-boundary section above shows.

Two further, more severe gaps were found and fixed together. First, mip-aware filtering:
EasyGL correctly honors any explicit \texttt{Mip}-qualified filter, but \texttt{Point} and
\texttt{Linear} on their own always sample mip level~0 regardless of minification --- a
\emph{deliberate}, documented deviation from FNA (which is mip-aware for these two filters
too), because CNA does not set \texttt{GL\_TEXTURE\_MAX\_LEVEL} per texture, and applying a
mip-aware filter unconditionally to the common non-mipmapped case would render as GL-incomplete
(solid black) instead. Second, and more severe: \texttt{Texture2D::SetData} at a mip level
above~0 was a \emph{total, silent no-op} on both Vulkan and BGFX --- the same historical
failure shape their \cnaclass{Texture3D} / \cnaclass{TextureCube} readback paths once had
(Chapter~\ref{ch:textures-rendertargets} now records the corrected live matrix). Mip-level
\texttt{SetData} is now a real GPU upload on all three hardware backends, but that upload fix
must not be mistaken for full \cnaclass{SamplerState} mip control: the common binding bridge has
no AddressW, MaxMipLevel, or LOD-bias argument. Backend caches consequently use their own
defaults for those missing values (for example, the Direct3D sampler caches set AddressW from
AddressV, LOD bias to zero, and allow the full mip range). No game-set value of those three
properties can reach them through this API today.
Anisotropic filtering itself: correct on Vulkan (queries the real device capability and
clamps the requested level); EasyGL was silently falling back to plain trilinear filtering
regardless of the requested anisotropy level until a fix wired real
OpenGL anisotropic-filtering-extension support, clamped to the driver's own capability;
BGFX remains partially fixed --- it enables the anisotropic sampling flag but
never communicates the actual requested level, an on/off switch rather than a graduated one,
still open. The newer backends widen that matrix: WebGPU and D3D11/D3D12 create real
anisotropic samplers (the Direct3D pair accept only feature level~11.0 or newer, whose maximum
anisotropy is 16); D3D9 maps the filter and maximum but neither capability-checks nor clamps
against \texttt{D3DCAPS9::MaxAnisotropy}; SDL GPU stores the requested maximum but never puts it
in a queued command or sampler-cache key; and Software ignores every sampler value after the
slot-range check. ASCII does not even forward \texttt{ApplySamplerState} to the SDL renderer it
wraps. The last three therefore also make
\texttt{SupportsCapability(AnisotropicFiltering)} definite false positives, since they inherit
the interface's default \texttt{true}; D3D9's answer is conditionally unreliable on the current
device. See Chapter~\ref{ch:backend-architecture}'s complete capability table.

\subsection{A worked example: crisp pixel art versus a minified, mipmapped floor}

Two everyday sampler needs pull \texttt{Filter} in opposite directions, and the chapter's own
mip-awareness deviation above means the obvious choice is not always the correct one for both.
A 2D pixel-art sprite wants hard, un-blurred texel edges at every zoom level --- the
\texttt{PointClamp} preset, unmodified, is exactly right, and mip level~0 is the \emph{only}
level a sprite like this has anyway:

\begin{lstlisting}
SamplerState pointClamp = SamplerState::PointClamp;
spriteBatch_->Begin(SpriteSortMode::Deferred, BlendState::AlphaBlend,
                     &pointClamp, nullptr, nullptr); // crisp texel edges, no mip chain to worry about
\end{lstlisting}

A distant, minified 3D floor texture with a real mipmap chain is the case where plain
\texttt{TextureFilter::Point} is the wrong choice in CNA specifically, precisely because of this
chapter's own already-stated deviation: \texttt{Point} and \texttt{Linear} both always sample mip
level~0 here, regardless of minification, so a floor texture with mips generated but sampled via
plain \texttt{Point} would show the same aliased, moir\'e-prone level~0 detail at any distance ---
mip selection never engages at all. Getting real mip-aware minification with a point-filtered
look (retro, blocky, but not shimmering into the distance) needs one of the explicitly
\texttt{Mip}-qualified filter values instead:

\begin{lstlisting}
SamplerState blockyButMipAware;
blockyButMipAware.setFilterProperty(TextureFilter::PointMipLinear); // point min/mag, real mip chain
blockyButMipAware.setAddressUProperty(TextureAddressMode::Wrap);
blockyButMipAware.setAddressVProperty(TextureAddressMode::Wrap);

device.getSamplerStatesProperty()[0] = blockyButMipAware;
\end{lstlisting}

This is a deliberate CNA deviation from FNA (which is mip-aware for plain \texttt{Point}/
\texttt{Linear} too), not an oversight left unfixed --- the chapter's own account of why explains
the alternative would be worse: unconditionally engaging mip selection for the common
non-mipmapped 2D case would render as GL-incomplete (solid black) rather than simply ignoring
mips, so the plain filters were deliberately kept mip-\emph{unaware} and the \texttt{Mip}-qualified
values are the one correct way to opt in for an asset that actually has mips.

\section{One recurring, small bug worth naming on its own}

Every one of these four state classes independently had the identical small bug at some
point in its history: its static presets did not set their own \texttt{Name} property. This
sounds trivial, but it had a real, non-cosmetic consequence for \cnaclass{DepthStencilState}
specifically --- \cnaclass{GraphicsDevice}'s own constructor did not actually \emph{copy}
\cnaclass{DepthStencilState::Default} and \cnaclass{RasterizerState::CullCounterClockwise}
into its own member fields at construction, and this had gone unnoticed because the values
happened to coincide by luck, until \texttt{Name} existed as a distinguishing field that
made the missing-copy bug observable at all. Fixing the missing-\texttt{Name} bug on all four
classes, in other words, is what \emph{surfaced} a real, separate device-construction bug ---
a small illustration of how a seemingly cosmetic gap in one place can be hiding a functional
one somewhere else entirely.
